\chapter*{\centering{\large{ABSTRAK}}}
\singlespacing{}

\textbf{Prabowo Darmawi}, Pengukuran Panjang Ikan  Menggunakan Metode \emph{Harris-Corner Detection} dan Estimasi Berat Ikan Menggunakan Metode Metode Turunan \emph{Length-Weight Relationship} Berbasis Citra Digital.
Skripsi, Program Studi Ilmu Komputer, Fakultas Matematika dan Ilmu Pengetahuan Alam, Universitas Negeri Jakarta, Mei 2026.
\\
\\
    Sektor budidaya perikanan di Indonesia seiring berkembangnya teknologi, masih memiliki permasalahan dalam pengukuran panjang dan berat ikan yang masih dilakukan secara manual.
Dalam penelitian ini, penulis menggunakan metode \emph{Harris-Corner Detection} untuk menghitung panjang dan metode turunan sederhana dari \emph{Length-Weight Relationship} untuk mengestimasi berat ikan berbasis Citra digital.
Langkah pertama yang dilakukan dalam Menghitung panjang dan estimasi berat ikan adalah Menginput gambar ikan, selanjutnya dengan menggunakan metode \emph{Harris-Corner Detection} untuk mendeteksi titik sudut pada citra ikan, 
dilanjutkan dengan menggunakan \emph{Bounding-Box} untuk memilih titik sudut dan menghitung panjang ikan, lalu menggunakan turunan sederhana dari \emph{Length-Weight Relationship} untuk mengestimasi berat ikan.
Hasil dari penelitian ini adalah sebuah algoritma yang dapat menghitung panjang dan berat ikan berdasarkan input gambar.
\\
\\
\textbf{Kata Kunci}: Corner Detection, Harris-Corner Detection, Bounding-Box, Length-Weight Relationship, Ikan.